\documentclass[11pt,a4paper]{article}
\usepackage[margin=21mm,headheight=17pt]{geometry}
\usepackage{fontspec}
\setmainfont[AutoFakeSlant=.18]{Noto Serif CJK SC}
\setsansfont{Noto Sans CJK SC}
\setmonofont{Noto Sans Mono CJK SC}
\XeTeXlinebreaklocale "zh"
\XeTeXlinebreakskip=0pt plus 1pt
\usepackage{amsmath,amssymb,amsthm,mathtools}
\usepackage{xcolor,enumitem,fancyhdr,longtable,booktabs,needspace,tikz}
\usepackage[unicode,colorlinks=true,linkcolor=blue!55!black,urlcolor=blue!55!black]{hyperref}
\newcommand{\R}{\mathbb{R}}\newcommand{\N}{\mathbb{N}}
\newcommand{\Q}{\mathbb{Q}}\newcommand{\Z}{\mathbb{Z}}
\newcommand{\eps}{\varepsilon}
\newcommand{\norm}[1]{\left\lVert#1\right\rVert}
\newcommand{\abs}[1]{\left\lvert#1\right\rvert}
\DeclareMathOperator{\conv}{co}\DeclareMathOperator{\cl}{cl}
\DeclareMathOperator{\Int}{int}\DeclareMathOperator{\diam}{diam}
\DeclareMathOperator{\dist}{dist}
\newtheorem{theorem}{定理}\newtheorem{lemma}{引理}
\newtheorem{proposition}{命题}
\theoremstyle{definition}\newtheorem{definition}{定义}\newtheorem{example}{例}
\renewcommand{\proofname}{证明}
\setlength{\parindent}{0pt}\setlength{\parskip}{.5em}\linespread{1.18}
\setlist{itemsep=3pt,topsep=4pt,leftmargin=2em}
\emergencystretch=4em\setcounter{tocdepth}{1}
\pagestyle{fancy}\fancyhf{}
\fancyhead[L]{\small 数理分析方法 · 证明考试自学讲义}
\fancyhead[R]{\small Lecture \LectureID}
\fancyfoot[C]{\thepage}
